% ============================================================
% EXAMEN DIAGNÓSTICO: FLUIDOS
% Curso de Oriente - CCH Oriente
% Enero 2026
% ============================================================
% INSTRUCCIONES: Pegar este contenido en tu plantilla de Overleaf
% después de la definición de entornos y antes de \end{document}
% ============================================================

\renewcommand{\tema}{Examen Diagnóstico --- Fluidos}

\twocolumn

\section*{Examen Diagnóstico: Fluidos}

% ============================================================
% PREGUNTAS
% ============================================================

\begin{preguntabox}
\subsection*{Pregunta 1}
Un buzo ejerce una fuerza de 500 N sobre una superficie de 0.05 m$^2$. ¿Cuál es la presión que ejerce?

\begin{enumerate}
    \item[A)] 25 Pa
    \item[B)] 2{,}500 Pa
    \item[C)] 10{,}000 Pa
    \item[D)] 50{,}000 Pa
\end{enumerate}
\end{preguntabox}

\begin{preguntabox}
\subsection*{Pregunta 2}
¿Cuál de las siguientes equivalencias para la presión atmosférica estándar es correcta?

\begin{enumerate}
    \item[A)] 1 atm $= 760$ Pa
    \item[B)] 1 atm $= 101{,}325$ Pa
    \item[C)] 1 atm $= 1{,}000$ mmHg
    \item[D)] 1 atm $= 10$ bar
\end{enumerate}
\end{preguntabox}

\begin{preguntabox}
\subsection*{Pregunta 3}
Un bloque de volumen 0.003 m$^3$ está completamente sumergido en agua ($\rho = 1{,}000$ kg/m$^3$). ¿Cuál es la fuerza de empuje?

\begin{enumerate}
    \item[A)] 3 N
    \item[B)] 15 N
    \item[C)] 30 N
    \item[D)] 60 N
\end{enumerate}
\end{preguntabox}

\begin{preguntabox}
\subsection*{Pregunta 4}
Una alberca tiene agua ($\rho = 1{,}000$ kg/m$^3$) a 4 m de profundidad. La presión sobre la superficie es $P_0 = 100{,}000$ Pa. ¿Cuál es la presión total en el fondo?

\begin{enumerate}
    \item[A)] 100{,}400 Pa
    \item[B)] 104{,}000 Pa
    \item[C)] 140{,}000 Pa
    \item[D)] 400{,}000 Pa
\end{enumerate}
\end{preguntabox}
\newpage
\begin{preguntabox}
\subsection*{Pregunta 5}
En una prensa hidráulica, $A_1 = 0.03$ m$^2$, $F_1 = 300$ N y $A_2 = 0.12$ m$^2$. ¿Cuál es la fuerza $F_2$ que produce el pistón grande?

\begin{enumerate}
    \item[A)] 75 N
    \item[B)] 300 N
    \item[C)] 900 N
    \item[D)] 1{,}200 N
\end{enumerate}
\end{preguntabox}

\begin{preguntabox}
\subsection*{Pregunta 6}
Una tubería de sección $A = 0.05$ m$^2$ conduce agua a $v = 6$ m/s. ¿Cuál es el gasto volumétrico?

\begin{enumerate}
    \item[A)] 0.03 m$^3$/s
    \item[B)] 0.11 m$^3$/s
    \item[C)] 0.30 m$^3$/s
    \item[D)] 3.00 m$^3$/s
\end{enumerate}
\end{preguntabox}

\begin{preguntabox}
\subsection*{Pregunta 7}
Un tronco de madera ($\rho = 600$ kg/m$^3$) flota en agua ($\rho = 1{,}000$ kg/m$^3$). ¿Qué fracción de su volumen queda bajo la superficie?

\begin{enumerate}
    \item[A)] 40\%
    \item[B)] 50\%
    \item[C)] 60\%
    \item[D)] 80\%
\end{enumerate}
\end{preguntabox}

\begin{preguntabox}
\subsection*{Pregunta 8}
Agua fluye por una tubería que se estrecha: $A_1 = 0.08$ m$^2$, $v_1 = 3$ m/s, $A_2 = 0.02$ m$^2$. ¿Cuál es la velocidad $v_2$ en la sección angosta?

\begin{enumerate}
    \item[A)] 0.75 m/s
    \item[B)] 6 m/s
    \item[C)] 12 m/s
    \item[D)] 24 m/s
\end{enumerate}
\end{preguntabox}
\newpage
\begin{preguntabox}
\subsection*{Pregunta 9}
Un tinaco abierto tiene agua a 45 m sobre un orificio en su base. ¿Con qué velocidad sale el chorro?

\begin{enumerate}
    \item[A)] 21 m/s
    \item[B)] 30 m/s
    \item[C)] 45 m/s
    \item[D)] 90 m/s
\end{enumerate}
\end{preguntabox}

\begin{preguntabox}
\subsection*{Pregunta 10}
En un tubo horizontal, agua ($\rho = 1{,}000$ kg/m$^3$) fluye con $v_1 = 2$ m/s donde $P_1 = 200{,}000$ Pa. En la sección angosta, $v_2 = 6$ m/s. ¿Cuál es la presión $P_2$?

\begin{enumerate}
    \item[A)] 182{,}000 Pa
    \item[B)] 184{,}000 Pa
    \item[C)] 200{,}000 Pa
    \item[D)] 216{,}000 Pa
\end{enumerate}
\end{preguntabox}

% ============================================================
% RESPUESTAS
% ============================================================

\newpage
\onecolumn

\section*{Respuestas y Retroalimentación}

\begin{respuestabox}
\subsection*{Pregunta 1}
\textcolor{verdecorrecto}{\textbf{Respuesta correcta: C) 10{,}000 Pa}}

\textbf{Justificación:}\\
\textbf{Fórmula utilizada:} Definición de presión.
$$P = \frac{F}{A} = \frac{500 \text{ N}}{0.05 \text{ m}^2} = 10{,}000 \text{ Pa}$$

\textbf{Respuestas incorrectas:}
\begin{itemize}
    \item \textbf{A) 25 Pa:} Multiplica $F \times A$ en lugar de dividir: $500 \times 0.05 = 25$.
    \item \textbf{B) 2{,}500 Pa:} Error de decimal: divide 500 entre 0.5 en lugar de 0.05.
    \item \textbf{D) 50{,}000 Pa:} Multiplica $F \times 100$ sin considerar el área correctamente.
\end{itemize}
\end{respuestabox}

\begin{respuestabox}
\subsection*{Pregunta 2}
\textcolor{verdecorrecto}{\textbf{Respuesta correcta: B) 1 atm $= 101{,}325$ Pa}}

\textbf{Justificación:}\\
La presión atmosférica estándar al nivel del mar es $P_0 = 101{,}325$ Pa $\approx 101$ kPa. Torricelli demostró que esta presión sostiene una columna de mercurio de 760 mm, de modo que 1 atm $= 760$ mmHg (no 760 Pa).

\textbf{Respuestas incorrectas:}
\begin{itemize}
    \item \textbf{A) 760 Pa:} Confunde el valor en mmHg (760) con el valor en Pa.
    \item \textbf{C) 1{,}000 mmHg:} El valor correcto en mmHg es 760, no 1{,}000.
    \item \textbf{D) 10 bar:} 1 atm $\approx 1.013$ bar, no 10 bar.
\end{itemize}
\end{respuestabox}

\begin{respuestabox}
\subsection*{Pregunta 3}
\textcolor{verdecorrecto}{\textbf{Respuesta correcta: C) 30 N}}

\textbf{Justificación:}\\
\textbf{Fórmula utilizada:} Principio de Arquímedes.
$$F_e = \rho_{\text{fluido}} \cdot g \cdot V_{\text{sumergido}} = 1{,}000 \frac{\text{kg}}{\text{m}^3} \times 10 \frac{\text{m}}{\text{s}^2} \times 0.003 \text{ m}^3 = 30 \text{ N}$$

\textbf{Respuestas incorrectas:}
\begin{itemize}
    \item \textbf{A) 3 N:} Olvida multiplicar por $g$; solo calcula $\rho \times V = 1{,}000 \times 0.003 = 3$.
    \item \textbf{B) 15 N:} Usa $g = 5$ m/s$^2$ en lugar de 10 m/s$^2$.
    \item \textbf{D) 60 N:} Duplica el resultado correcto sin justificación física.
\end{itemize}
\end{respuestabox}

\begin{respuestabox}
\subsection*{Pregunta 4}
\textcolor{verdecorrecto}{\textbf{Respuesta correcta: C) 140{,}000 Pa}}

\textbf{Justificación:}\\
\textbf{Fórmula utilizada:} Presión hidrostática.
$$P = P_0 + \rho g h = 100{,}000 \text{ Pa} + \left(1{,}000 \frac{\text{kg}}{\text{m}^3}\right)\left(10 \frac{\text{m}}{\text{s}^2}\right)(4 \text{ m}) = 100{,}000 + 40{,}000 = 140{,}000 \text{ Pa}$$

\textbf{Respuestas incorrectas:}
\begin{itemize}
    \item \textbf{A) 100{,}400 Pa:} Calcula $\rho \times h = 4{,}000$ sin multiplicar por $g$, luego suma mal.
    \item \textbf{B) 104{,}000 Pa:} Suma solo 4{,}000 en lugar de 40{,}000 (error de decimal en $\rho g h$).
    \item \textbf{D) 400{,}000 Pa:} Multiplica $P_0 \times h$ en lugar de sumar $\rho g h$.
\end{itemize}
\end{respuestabox}

\begin{respuestabox}
\subsection*{Pregunta 5}
\textcolor{verdecorrecto}{\textbf{Respuesta correcta: D) 1{,}200 N}}

\textbf{Justificación:}\\
\textbf{Fórmula utilizada:} Principio de Pascal --- prensa hidráulica.
$$F_2 = F_1 \cdot \frac{A_2}{A_1} = 300 \text{ N} \times \frac{0.12 \text{ m}^2}{0.03 \text{ m}^2} = 300 \text{ N} \times 4 = 1{,}200 \text{ N}$$

\textbf{Respuestas incorrectas:}
\begin{itemize}
    \item \textbf{A) 75 N:} Invierte la razón de áreas: usa $A_1/A_2 = 0.03/0.12 = 0.25$.
    \item \textbf{B) 300 N:} Supone que la fuerza no cambia al cambiar el área, ignorando Pascal.
    \item \textbf{C) 900 N:} Usa la diferencia de áreas como factor en lugar del cociente.
\end{itemize}
\end{respuestabox}

\begin{respuestabox}
\subsection*{Pregunta 6}
\textcolor{verdecorrecto}{\textbf{Respuesta correcta: C) 0.30 m$^3$/s}}

\textbf{Justificación:}\\
\textbf{Fórmula utilizada:} Gasto volumétrico.
$$Q = A \cdot v = 0.05 \text{ m}^2 \times 6 \frac{\text{m}}{\text{s}} = 0.30 \frac{\text{m}^3}{\text{s}}$$

\textbf{Respuestas incorrectas:}
\begin{itemize}
    \item \textbf{A) 0.03 m$^3$/s:} Error de decimal: opera como si $A = 0.005$ m$^2$.
    \item \textbf{B) 0.11 m$^3$/s:} Suma en lugar de multiplicar: $0.05 + 6$ y luego divide incorrectamente.
    \item \textbf{D) 3.00 m$^3$/s:} Confunde el área: usa $A = 0.5$ m$^2$ en lugar de 0.05 m$^2$.
\end{itemize}
\end{respuestabox}

\begin{respuestabox}
\subsection*{Pregunta 7}
\textcolor{verdecorrecto}{\textbf{Respuesta correcta: C) 60\%}}

\textbf{Justificación:}\\
\textbf{Fórmula utilizada:} Fracción sumergida en equilibrio de flotación.
$$\frac{V_{\text{sum}}}{V_{\text{total}}} = \frac{\rho_{\text{obj}}}{\rho_{\text{fluido}}} = \frac{600 \text{ kg/m}^3}{1{,}000 \text{ kg/m}^3} = 0.60 = 60\%$$

\textbf{Respuestas incorrectas:}
\begin{itemize}
    \item \textbf{A) 40\%:} Calcula la fracción que queda \textit{fuera} del agua: $1 - 0.60 = 0.40$.
    \item \textbf{B) 50\%:} Supone que todo objeto flotante se sumerge a la mitad, sin usar densidades.
    \item \textbf{D) 80\%:} Invierte la razón de densidades o confunde los valores entre sí.
\end{itemize}
\end{respuestabox}

\begin{respuestabox}
\subsection*{Pregunta 8}
\textcolor{verdecorrecto}{\textbf{Respuesta correcta: C) 12 m/s}}

\textbf{Justificación:}\\
\textbf{Fórmula utilizada:} Ecuación de continuidad.
$$v_2 = \frac{A_1 v_1}{A_2} = \frac{0.08 \text{ m}^2 \times 3 \text{ m/s}}{0.02 \text{ m}^2} = \frac{0.24 \text{ m}^3/\text{s}}{0.02 \text{ m}^2} = 12 \text{ m/s}$$

\textbf{Respuestas incorrectas:}
\begin{itemize}
    \item \textbf{A) 0.75 m/s:} Invierte la razón: $v_2 = v_1 \times (A_2/A_1) = 3 \times 0.25 = 0.75$.
    \item \textbf{B) 6 m/s:} Solo duplica $v_1$, usando factor 2 en lugar del factor correcto (4).
    \item \textbf{D) 24 m/s:} Cuadruplica incorrectamente: aplica el factor de áreas dos veces.
\end{itemize}
\end{respuestabox}

\begin{respuestabox}
\subsection*{Pregunta 9}
\textcolor{verdecorrecto}{\textbf{Respuesta correcta: B) 30 m/s}}

\textbf{Justificación:}\\
\textbf{Fórmula utilizada:} Teorema de Torricelli.
$$v = \sqrt{2gh} = \sqrt{2 \times 10 \frac{\text{m}}{\text{s}^2} \times 45 \text{ m}} = \sqrt{900 \text{ m}^2/\text{s}^2} = 30 \text{ m/s}$$

\textbf{Respuestas incorrectas:}
\begin{itemize}
    \item \textbf{A) 21 m/s:} Olvida el factor 2: calcula $\sqrt{gh} = \sqrt{450} \approx 21$ m/s.
    \item \textbf{C) 45 m/s:} Confunde la velocidad con la altura; responde directamente con $h$.
    \item \textbf{D) 90 m/s:} Calcula $2gh = 900$ correctamente pero no extrae la raíz cuadrada.
\end{itemize}
\end{respuestabox}

\begin{respuestabox}
\subsection*{Pregunta 10}
\textcolor{verdecorrecto}{\textbf{Respuesta correcta: B) 184{,}000 Pa}}

\textbf{Justificación:}\\
\textbf{Fórmula utilizada:} Ecuación de Bernoulli en tubo horizontal ($h_1 = h_2$).
$$P_2 = P_1 + \frac{1}{2}\rho\left(v_1^2 - v_2^2\right)$$
$$P_2 = 200{,}000 \text{ Pa} + \frac{1}{2}\left(1{,}000 \frac{\text{kg}}{\text{m}^3}\right)\left[(2 \text{ m/s})^2 - (6 \text{ m/s})^2\right]$$
$$P_2 = 200{,}000 \text{ Pa} + 500 \times (4 - 36) \text{ Pa} = 200{,}000 - 16{,}000 = 184{,}000 \text{ Pa}$$

\textbf{Respuestas incorrectas:}
\begin{itemize}
    \item \textbf{A) 182{,}000 Pa:} Solo usa $\tfrac{1}{2}\rho v_2^2 = 18{,}000$ sin restar la contribución de $v_1^2$.
    \item \textbf{C) 200{,}000 Pa:} Supone que la presión no cambia al cambiar la velocidad, ignorando Bernoulli.
    \item \textbf{D) 216{,}000 Pa:} Suma en lugar de restar: usa $(v_2^2 - v_1^2)$ en lugar de $(v_1^2 - v_2^2)$.
\end{itemize}
\end{respuestabox}
